\section{第\ref{sec:pain}章　痛み}

痛みのセンサ、2本の配線、脊髄のゲート、下行性の音量つまみ、感作、注意の奪取は直接測定されている。ゲートと感作の式は本書の単純化したモデルである。マシンが警報の音量を選ぶ規則は仮説である。

{\footnotesize
\setlength{\tabcolsep}{3pt}\renewcommand{\arraystretch}{1.15}
\begin{longtable}{>{\raggedright}p{0.5cm}>{\raggedright}p{4.0cm}>{\raggedright}p{5.6cm}>{\raggedright}p{2.3cm}>{\raggedright\arraybackslash}p{1.7cm}}
\toprule
No. & 主張 & 実験と測定されたもの & 出典 & 状態 \\
\midrule\endfirsthead
\toprule
No. & 主張 & 実験と測定されたもの & 出典 & 状態 \\
\midrule\endhead
1 & 高い閾値：痛みのセンサの加熱閾値は $41^\circ$ から $46$--$48^\circ$ & 単一線維記録。サルの皮膚：遅い（C）センサの閾値の中央値は $41^\circ$、速いセンサの第2型は $46^\circ$、第1型は $53^\circ$ 超。ヒトの皮膚：圧力にも応答するCセンサは $40.7^\circ$、圧力に応答しないものは $48^\circ$。 & \cite{Treede1995heat, Weidner1999cfib} & 測定済み \\
2 & 痛みのセンサは触覚のセンサよりはるかに順応しにくく、軽い損傷の後にはより敏感になる；強い加熱の後の低下は一つの型で測定 & 皮膚の軽いやけど（$50^\circ$、$100\,\text{s}$）：$5$--$10$~分後、ヒトの痛みの閾値とサルのCセンサの閾値が $2$--$6^\circ$ 下がる；他の皮膚センサにはこのずれはない。速いセンサの第2型では、強い加熱の後に応答が抑えられる。触覚センサとの順応の比較は一般的な推定で、ここでは個別の実験で裏づけられていない。 & \cite{LaMotte1982hyper, Treede1995heat} & 測定済み \\
3 & 2本の配線：速いほうは $5$--$30$~m/s、遅いほうは $0.5$--$2$~m/s；到着時間 $t=L/v$~\eqref{eq:paintime} & 伝導速度：サルの速い痛みのセンサは平均 $15$ と $25$~m/s（二つの型）；ヒトのCセンサは $0.8$--$1.0$~m/s。$L=1$~m ならそれぞれ数十ミリ秒と約1秒。 & \cite{Treede1995heat, Weidner1999cfib} & 測定済み \\
4 & 痛みは2回来る：最初は速く正確に、次に鈍く長く & ヒト前腕の加熱に対する反応時間：温度の上昇とともに $1100$ から $400\ms$；有毛部の皮膚の強い加熱は二重の痛みを生むが、手のひらでは生まない。第一の痛みを運べるのは $6$~m/s より速い線維だけで、潜時からそれに対応するのは手のひらにない速いセンサの第2型である。役割の分担（「どこ」と「どれほど悪いか」）は解釈。 & \cite{CampbellLaMotte1983first, Treede1995heat} & 測定済み \\
5 & ゲート：脊髄の出力ニューロンは痛みと触覚を受け、抑制性の介在ニューロンは触覚で興奮する（図~\ref{fig:gate}） & ゲートの構成は理論として提案された。マウス、遺伝的標識とニューロン群の除去：ソマトスタチンを持つ興奮性ニューロンは機械的な痛みに必要；ダイノルフィンを持つ抑制性ニューロンは触覚センサからの入力がそこへ届くのを防ぐ。これらのニューロンがないと軽い接触が痛みを起こす。 & \cite{MelzackWall1965, Duan2014} & 測定済み \\
6 & 触覚は痛みを弱める & ヒト、腕へのレーザー痛み刺激：同時の触覚は痛みの評定とパルスのエネルギーの識別能力を下げる；効果は、同じ皮膚分節の中で触れる点と痛む点の距離とともに弱まる。 & \cite{Mancini2014touch} & 測定済み \\
7 & ゲート理論の仮定：強い痛みは自ら抑制性介在ニューロンをオフにし、ゲートが開け放たれる & 章では元の理論の仮定と明記。マウスでの研究は、ゲートが触覚を痛みのチャネルに入れない仕組みを記述するが、痛みによる介在ニューロンのオフは示していない。 & \cite{MelzackWall1965} & 仮説 \\
8 & ゲート~\eqref{eq:gate}：閾値は触覚なしで $0.18$、触覚ありで $0.86$、$g=2$ で $0.11$；$\nu=1$ で触覚は出力を $1$ から $0.25$ に下げ、$\nu=2$ では効かない；触覚単独では通らない（図~\ref{fig:gatecurves}） & 本文の重みでの \texttt{models/\allowbreak pain\_\allowbreak gate.py} による計算が、これらの数値をすべて再現する。 & \texttt{models/\allowbreak pain\_\allowbreak gate.py} & モデル \\
9 & 脳幹からの下行路は、ゲートのゲインをどちらの向きにも変える & ラット延髄、尾の加熱中の記録：あるニューロンは引っ込めの前に発火し（「オン」）、別のニューロンは沈黙する（「オフ」）；この領域の弱い刺激は反射を抑える。総説：同じ回路が痛みの伝達を促進も抑制もし、オピオイドはそれを介して働く。 & \cite{Fields1983rvm, Fields2004} & 測定済み \\
10 & 楽になるという期待は、オピオイドのチャネルを通じて痛みを弱める & 急性の痛みを持つ人：楽になるという期待で痛みが減った人では、オピオイド受容体の遮断で痛みが他の人と同じレベルに戻った；それ以外の人では遮断は痛みを変えなかった。 & \cite{Levine1978} & 測定済み \\
11 & 脳は割り込みの損得によって警報の音量を選ぶ & 音量を選ぶ規則を測定した実験はない。 & --- & 仮説 \\
12 & 感作：損傷の後、ゲートの出力は増し、閾値は下がり、その一部は脊髄による；損傷刺激が止んだ後も続く & ラット：皮膚の損傷の後、屈曲反射の閾値が下がり、応答が増す；電気生理学的解析によれば、増大の一部は脊髄そのもので生じる。損傷刺激は刺激そのものより長く続く感度の変化を起こす。正確な回復時間は章では挙げていない。 & \cite{Woolf1983} & 測定済み \\
13 & 感作~\eqref{eq:sensitization}は正帰還である；ループが強いと、警報は治癒後もラッチしうる & モデル~\eqref{eq:sensitization}から導かれる（章末の問題）。治癒後も続く痛みがまさにこの形のラッチしたループであることを示す実験はない。 & 章内の導出 & 仮説 \\
14 & 痛みは負の報酬であり、それによる学習は時間差分誤差に従う & fMRI、ヒト、痛みの前触れの連鎖：腹側線条体と前部島皮質の活動が時間差分学習の信号と一致する。サル：\nt{DOPA}ニューロンの一部は不快なエアパフとその前触れで興奮し、別の一部は抑制される。 & \cite{Seymour2004, Matsumoto2009} & 測定済み \\
15 & 痛みはいまの作業に割り込む & ヒト、電気的な痛み刺激と音のプローブ：痛みの開始 $100\ms$ 後のプローブへの応答は明らかに遅れ、$1.5\,\text{s}$ 後には対照との差はもうない。総説はこうした実験を注意の奪取のモデルにまとめる。 & \cite{Crombez1996disrupt, EcclestonCrombez1999} & 測定済み \\
16 & 引っ込め反射は脊髄で閉じる & ラット：後角深層のニューロンは個々の筋肉の引っ込め反射の空間パターンを再現する；頸部からは逆行性に興奮させられず、脊髄の介在ニューロンである。 & \cite{Schouenborg1995wr} & 測定済み \\
\bottomrule
\end{longtable}}
